\documentclass[11pt]{article}

\usepackage[margin=1in]{geometry}
\usepackage[T1]{fontenc}
\usepackage[utf8]{inputenc}
\usepackage{microtype}
\usepackage{booktabs}
\usepackage{longtable}
\usepackage{array}
\usepackage{hyperref}
\usepackage[numbers,sort&compress]{natbib}

\title{Workplace Determinants of Metabolic Syndrome in Working Adults: Systematic Review}
\author{Gideon Ofosu Kyere\\
Kwame Nkrumah University of Science and Technology (KNUST), Kumasi, Ghana\\
ORCID: 0009-0003-9848-8437}
\date{}

\begin{document}
\maketitle

\begin{abstract}
\textbf{Objectives:} To synthesise evidence on occupational stress, diet, physical activity, sedentary behaviour and workplace lifestyle interventions in relation to metabolic syndrome (MetS) among working adults.

\textbf{Methods:} Peer-reviewed studies published from 2015 to 2026 were identified through a workplace-focused PubMed/MEDLINE search supplemented by targeted searches, reference screening and citation chasing. Eligible studies involved employed adults or occupational populations and examined workplace stress, diet, physical activity, sedentary behaviour or workplace lifestyle interventions in relation to MetS or established MetS components. Randomised trials, cohort studies, analytical cross-sectional studies and quasi-experimental interventions were eligible. Design-specific Joanna Briggs Institute tools were used for critical appraisal. Because of substantial clinical and methodological heterogeneity, findings were synthesised narratively.

\textbf{Results:} Thirty-four primary studies contributed to the narrative synthesis. Prospective studies linked high or increasing work stress with subsequent MetS risk, whereas cross-sectional stress findings were less consistent. Higher leisure-time physical activity was generally associated with lower MetS risk, while prolonged sedentary exposure was associated with higher risk. Healthier dietary patterns were associated with more favourable metabolic profiles. Workplace interventions involving exercise, dietary modification, counselling or multicomponent lifestyle support improved waist circumference, blood pressure, glucose regulation, lipid measures or MetS severity in several studies.

\textbf{Conclusions:} MetS in working adults is associated with a combination of occupational and behavioural factors. The strongest evidence supports adverse work stress, physical inactivity and prolonged sedentary behaviour as relevant risk markers, while workplace lifestyle interventions can improve several metabolic outcomes. More consistent longitudinal and randomised evidence is needed across diverse occupational settings.
\end{abstract}

\section*{What is already known on this topic}
Occupational stress, diet, physical inactivity and sedentary behaviour have each been linked to cardiometabolic risk in workers, but the evidence is dispersed across occupational groups and study designs.

\section*{What this study adds}
This review integrates prospective, cross-sectional and workplace-intervention evidence across four domains: occupational stress, diet/nutrition, physical activity/sedentary behaviour and multicomponent workplace programmes. Prospective evidence was strongest for work stress, while intervention studies showed that workplace-based lifestyle modification can improve several MetS components.

\section*{How this study might affect research, practice or policy}
Workplace prevention may benefit from combining psychosocial risk management, healthier food environments, opportunities for physical activity, reduced prolonged sitting and targeted cardiometabolic screening. Future research should prioritise standardised MetS definitions and stronger longitudinal and randomised designs.

\section{Introduction}

Metabolic syndrome (MetS) brings together a cluster of cardiometabolic abnormalities---typically central obesity, elevated blood pressure, impaired glucose regulation, raised triglycerides and reduced high-density lipoprotein cholesterol---that substantially increase the risk of type 2 diabetes and cardiovascular disease.

The workplace is an important setting in which these risks accumulate. Many adults spend a large proportion of waking hours at work, where physical demands, sitting time, shift patterns, psychosocial pressure and access to food can shape everyday behaviour. A sedentary job may reduce energy expenditure; irregular or extended hours may disrupt eating and recovery; and sustained psychosocial stress may influence both behaviour and physiological regulation. At the same time, the workplace offers a practical platform for prevention through screening, food-environment changes, exercise programmes, counselling and digital health support.

The occupational evidence is heterogeneous. Prospective studies have linked high or increasing work stress with later MetS \citep{Garbarino2015,Yamaguchi2018}, while large workforce cohorts have shown differences in MetS risk across occupational groups \citep{vanZon2020Occupational,Runge2021Metabolic}. Leisure-time exercise appears protective in some worker cohorts \citep{Kuwahara2016}, whereas prolonged sedentary exposure has been associated with higher MetS prevalence \citep{Motuma2020}. Intervention studies further suggest that workplace-based exercise, dietary modification and multicomponent lifestyle programmes can improve metabolic outcomes \citep{Haufe2019,Ma2025,Bogale2025}.

What is less clear is how these strands of evidence fit together. Studies differ in occupation, exposure definition, intervention intensity, MetS criteria and methodological design, making a simple pooled estimate inappropriate. This review therefore synthesised the evidence across occupational stress, diet and workplace nutrition, physical activity and sedentary behaviour, and workplace lifestyle interventions among working adults.

\section{Methods}

\subsection{Review design and eligibility}
A workplace-focused systematic review was undertaken using revised eligibility criteria and a fresh literature search. Adults in employment, occupational groups or explicitly defined workplace settings were eligible. Exposures and interventions included occupational stress or job strain, dietary patterns and nutrition, physical activity or exercise, sedentary behaviour, and workplace lifestyle interventions.

Studies had to report MetS as a defined syndrome or report established MetS components in a context explicitly concerned with MetS risk, prevention, severity or management. Eligible designs were randomised controlled trials, prospective or longitudinal cohort studies, analytical cross-sectional studies and quasi-experimental workplace interventions. Protocols, editorials, conference abstracts, grey literature and secondary reviews were excluded. The publication window was 2015--2026.

\subsection{Search and study selection}
PubMed/MEDLINE was the primary bibliographic source. Search concepts combined MetS with terms for workers or workplaces and occupational stress, physical activity, sedentary behaviour, diet, nutrition and lifestyle intervention. Supplementary searches used narrower exposure-specific combinations, together with reference-list screening, citation chasing, publisher records and related-article searching. The complete search strategy is provided in Supplementary Methods S1.

Screening was conducted at title/abstract and full-text stages. Multiple reports from the same cohort or intervention were linked so that overlapping populations were not treated as independent evidence.

% FINALISATION NOTE: insert the completed PRISMA flow counts here before submission.

\subsection{Data extraction and critical appraisal}
A structured extraction framework captured occupational population, design, sample size, exposure or intervention, MetS outcome, principal effect estimate and methodological limitations. Adjusted estimates were preferred for observational studies, and between-group effects were prioritised for interventions where available.

Methodological appraisal used design-specific Joanna Briggs Institute tools for randomised trials, cohort studies, analytical cross-sectional studies and quasi-experimental studies. Appraisal findings were used to qualify interpretation rather than converted into a numerical quality score. Item-level results are reported in Supplementary Tables S2--S5.

\subsection{Synthesis}
Substantial heterogeneity in populations, exposures, interventions, MetS definitions and effect measures precluded a meaningful overall meta-analysis. Evidence was therefore synthesised narratively by domain, with greater interpretive weight given to randomised and prospective evidence than to cross-sectional or uncontrolled findings.

Formal publication-bias analysis was not undertaken because no pooled meta-analysis was performed. A formal GRADE assessment was also not applied; confidence was judged in relation to study design, directness, consistency and JBI appraisal findings.

\section{Results}

\subsection{Evidence base}

Thirty-four primary studies published between 2015 and 2026 contributed to the synthesis. Six were randomised trials, six were prospective or longitudinal cohorts, 17 were analytical cross-sectional studies and five were quasi- or non-randomised interventions. The evidence covered a broad range of occupational settings, including office and white-collar work, police and security services, commercial driving, manufacturing and petrochemical work, healthcare, universities, the military, heavy industry and employer-based health-promotion programmes.

The studies also differed substantially in what they treated as the principal exposure. Some focused on psychosocial work stress, others on diet or physical activity, and several evaluated more than one domain simultaneously. Intervention studies ranged from single-component exercise or dietary programmes to multicomponent strategies incorporating counselling, self-monitoring, digital support or broader lifestyle education. This heterogeneity supported narrative synthesis rather than a single pooled estimate.

\section*{Table 1. Overview of the evidence base}

\begin{tabular}{p{0.24\textwidth} p{0.28\textwidth} p{0.40\textwidth}}
\toprule
Evidence domain & Main study designs & Overall pattern \\
\midrule
Occupational stress & Prospective cohort; analytical cross-sectional & Prospective studies linked high or increasing work stress with later MetS risk; cross-sectional findings were more mixed and often component-specific. \\
Diet and workplace nutrition & Analytical cross-sectional; randomised and quasi-experimental intervention & Healthier dietary patterns generally aligned with more favourable metabolic profiles, while workplace dietary interventions improved several MetS components. \\
Physical activity and sedentary behaviour & Cohort; analytical cross-sectional; randomised intervention & Leisure-time physical activity was generally associated with lower MetS risk, whereas prolonged sedentary exposure was associated with higher risk. \\
Multicomponent workplace interventions & Randomised and quasi-experimental intervention & Programmes combining diet, activity, counselling, self-monitoring or digital support generally improved anthropometric or metabolic outcomes. \\
\bottomrule
\end{tabular}

\vspace{0.5em}
\noindent\textit{Note:} Several studies contributed evidence to more than one domain. Detailed characteristics for all 34 studies are provided in Supplementary Table S1.
\section*{Table 2. Selected source-verified findings}

\begin{longtable}{p{0.18\textwidth} p{0.25\textwidth} p{0.48\textwidth}}
\toprule
Study & Exposure/intervention & Principal finding \\
\midrule
\endfirsthead
\toprule
Study & Exposure/intervention & Principal finding \\
\midrule
\endhead
Garbarino \& Magnavita (2015) & Work stress & High work stress was associated with incident MetS (adjusted OR 2.68, 95\% CI 1.08--6.70). \\
Yamaguchi et al. (2018) & Change in job demands & Low-to-high job demands were associated with incident MetS (adjusted OR 3.27, 95\% CI 1.46--7.34). \\
Kuwahara et al. (2016) & Leisure-time exercise & Compared with the lowest category, mixed moderate/vigorous exercise was associated with lower MetS incidence (HR 0.74, 95\% CI 0.62--0.89 for 7.5--<16.5 MET-h/week). \\
Browne et al. (2017) & Physical activity & Active workers had lower odds of MetS than sedentary workers (OR 0.52, 95\% CI 0.27--0.98). \\
Motuma et al. (2020) & Sedentary behaviour and diet & Sedentary time $\geq$8 h/day was associated with higher MetS prevalence (APR 2.29, 95\% CI 1.88--2.80); high fruit/vegetable intake showed an inverse association. \\
Haufe et al. (2019) & Telemonitored exercise & Exercise reduced MetS severity versus control (between-group difference $-0.26$, 95\% CI $-0.35$ to $-0.16$). \\
Ma et al. (2025) & Canteen-based dietary intervention & The intervention improved fasting glucose, total and LDL cholesterol, waist circumference, BMI and HDL cholesterol. \\
Bogale et al. (2025) & Lifestyle education & Waist circumference decreased by 5.33 cm and systolic blood pressure by 6.96 mmHg, with additional lipid improvements. \\
Ayeltigah et al. (2026) & Work stress in military personnel & Work stress was associated with higher odds of MetS (OR 6.472, 95\% CI 1.535--27.292), although the estimate was imprecise. \\
Ryu et al. (2017) & Targeted workplace health promotion & Waist circumference and fasting glucose improved, while MetS prevalence declined from 41.5\% to 31.7\%. \\
\bottomrule
\end{longtable}

\vspace{0.5em}
\noindent\textit{Abbreviations:} APR, adjusted prevalence ratio; BMI, body mass index; CI, confidence interval; HR, hazard ratio; MetS, metabolic syndrome; OR, odds ratio.

\section*{Table 3. Methodological appraisal by design}

\begin{tabular}{p{0.24\textwidth} p{0.08\textwidth} p{0.24\textwidth} p{0.36\textwidth}}
\toprule
Design & n & Main strengths & Main limitations \\
\midrule
Randomised controlled trials & 6 & Stronger causal inference; standardised outcome assessment & Allocation concealment, masking, attrition and analysis by randomised group were incompletely reported in several trials. \\
Prospective/cohort studies & 6 & Temporal ordering; adjustment for major confounders; standardised MetS assessment & Incomplete reporting of attrition or missing-data handling in several studies. \\
Analytical cross-sectional studies & 17 & Broad occupational coverage; generally standardised MetS criteria & No temporal ordering; reverse causation remains plausible. \\
Quasi-/non-randomised interventions & 5 & Real-world workplace implementation & Non-randomised allocation and, in several studies, absence of a concurrent control group. \\
\bottomrule
\end{tabular}

\vspace{0.5em}
\noindent\textit{Note:} Design-specific Joanna Briggs Institute appraisal was used to qualify confidence in the narrative synthesis; no summed quality score was applied.

\subsection{Occupational stress}

The clearest temporal evidence came from prospective cohorts. Among police officers, high work stress was associated with incident MetS over five years (adjusted OR 2.68, 95\% CI 1.08--6.70) \citep{Garbarino2015}. In Japanese manufacturing workers, moving from low to high job demands was associated with a greater likelihood of developing MetS over three years (adjusted OR 3.27, 95\% CI 1.46--7.34) \citep{Yamaguchi2018}. These studies provide stronger support for temporality than the cross-sectional evidence because the occupational exposure preceded subsequent metabolic outcomes.

Cross-sectional findings were more mixed. Security-guard data linked occupational stressors with adverse glucose, lipid and blood-pressure profiles \citep{Jovanovic2020}, while medical-university staff showed variation in MetS prevalence across employment categories \citep{Eftekhari2021}. Among Ghanaian military personnel, work stress was associated with substantially higher odds of MetS, although the confidence interval was wide and the estimate therefore imprecise \citep{Ayeltigah2026Military}. By contrast, Zhang et al. found no overall adjusted association between occupational stress and MetS in petrochemical workers; significant findings were limited to selected stress dimensions and individual MetS components \citep{Zhang2024}. In police officers followed prospectively, cardiorespiratory fitness predicted lower subsequent MetS risk, whereas work stress itself did not show a significant main effect \citep{Schilling2020}. Overall, the stress literature was directionally suggestive but sensitive to the construct measured, the occupational setting and the study design.

\subsection{Diet and workplace nutrition}

Dietary evidence came from both observational and intervention studies. In office workers, lower fruit intake was associated with MetS \citep{Alavi2015}. Among Ghanaian tanker drivers, MetS was common and was observed alongside suboptimal dietary intake patterns \citep{Apprey2022}. Among Iranian health workers, poorer diet moderation was associated with MetS and several components \citep{Sharifi2025}. In a large Madrid workforce, greater adherence to a Mediterranean dietary pattern was associated with substantially lower odds of MetS \citep{ReinosoBarbero2026}.

Other studies suggested that meal timing and occupational schedules may also matter. Among healthcare workers, night work was accompanied by less favourable metabolic profiles, while a shorter eating window among night workers was associated with lower MetS risk \citep{dosReis2026}. In Ghanaian healthcare and university employees, dietary-pattern associations were more apparent than associations with stress or physical activity for some cardiometabolic markers \citep{Ayeltigah2026Healthcare,Ayeltigah2026University}. Heavy-industry workers with obesity or abdominal obesity also had markedly higher odds of MetS \citep{Petek2026}.

Intervention evidence provided stronger support for modifiability. A randomised canteen-based programme improved fasting glucose, total and LDL cholesterol, HDL cholesterol, waist circumference and body mass index \citep{Ma2025}. A Korean worksite dietary programme similarly improved several MetS indicators among higher-risk workers \citep{Kim2016Dietary}. These findings indicate that changing the workplace food environment can affect metabolic risk factors, although the observational dietary literature remains vulnerable to self-report error and residual confounding.

\subsection{Physical activity and sedentary behaviour}

Physical activity showed one of the most consistent patterns across designs. In a prospective cohort of more than 22,000 workers, moderate-to-vigorous leisure-time exercise was associated with lower MetS incidence \citep{Kuwahara2016}. Among workers in sedentary occupations, meeting physical-activity recommendations was associated with lower odds of MetS (OR 0.52, 95\% CI 0.27--0.98) and several individual components \citep{Browne2017}. Taiwanese workers with high leisure-time activity likewise had lower odds of MetS, while commuting activity showed no clear association \citep{Huang2017Taiwan}. Occupational and commuting activity did not consistently reproduce the protective pattern seen for leisure-time exercise.

Sedentary exposure showed the opposite direction. Among working adults in Ethiopia, sitting for at least eight hours per day was associated with more than twice the prevalence of MetS (APR 2.29, 95\% CI 1.88--2.80) \citep{Motuma2020}. Differences across occupational groups in large Dutch cohorts further suggested that the metabolic consequences of work vary with the broader physical and socioeconomic characteristics of occupations \citep{vanZon2020Occupational,Runge2021Metabolic}.

Intervention evidence reinforced the potential importance of activity. Telemonitoring-supported exercise reduced MetS severity in company employees \citep{Haufe2019}. T'ai Chi improved blood pressure and triglycerides in middle-aged office workers \citep{Choi2017}; a structured physical-activity programme in overweight high-tech employees reduced the number of MetS risk factors and improved lipid outcomes \citep{Fang2019}; and a remote activity programme improved physiological indicators among male office workers with MetS \citep{ParkHwang2024}.

\subsection{Multicomponent workplace interventions}

Interventions combining diet, exercise, counselling, self-monitoring or digital support generally moved metabolic outcomes in a favourable direction. A voluntary worksite weight-loss programme reduced MetS prevalence in both women and men \citep{Earnest2015}. A targeted workplace health-promotion programme reduced waist circumference and fasting glucose and lowered MetS prevalence from 41.5\% to 31.7\% \citep{Ryu2017}. An internet-based lifestyle programme improved selected cardiometabolic risk factors \citep{Kim2015Internet}, and a web-based high-intensity interval-training programme combined with calorie restriction produced sustained improvements in weight, blood pressure and triglycerides \citep{So2025}.

The strongest recent intervention evidence came from randomised programmes. In Ethiopian bank employees, a cluster-randomised lifestyle-education intervention reduced waist circumference by 5.33 cm and systolic blood pressure by 6.96 mmHg, alongside favourable lipid changes \citep{Bogale2025}. Together with the exercise and dietary trials, these studies suggest that workplace metabolic risk is modifiable. However, programme content, intensity, duration and follow-up varied considerably, and several non-randomised interventions lacked a concurrent control group. Direct comparison of effect magnitude across programmes is therefore inappropriate.

\section{Discussion}

This review indicates that metabolic risk in working adults is shaped by an interplay between occupational exposures and modifiable behaviours rather than by a single workplace factor. Across the evidence base, three findings were particularly consistent. First, adverse work stress showed its strongest signal when examined prospectively. Second, greater leisure-time physical activity and lower sedentary exposure were associated with more favourable metabolic profiles across observational and intervention studies. Third, workplace lifestyle programmes could improve several MetS components, particularly when they addressed more than one behavioural pathway.

The work-stress findings require nuance. The prospective police and manufacturing cohorts provide meaningful temporal support because work stress preceded incident MetS \citep{Garbarino2015,Yamaguchi2018}. Yet the cross-sectional literature was less uniform, and Zhang et al. found no overall adjusted association between occupational stress and MetS \citep{Zhang2024}. This inconsistency argues against treating work stress as a single homogeneous exposure. Job demands, effort--reward imbalance, control, social support, shift structure and recovery opportunities may operate through different pathways, and their metabolic consequences may depend on sleep, diet, physical activity and baseline cardiometabolic vulnerability. The relatively large estimate among Ghanaian military personnel is therefore notable but should be interpreted cautiously given its wide confidence interval and cross-sectional design \citep{Ayeltigah2026Military}.

The physical-activity evidence was more coherent. Prospective, cross-sectional and intervention studies all pointed towards a beneficial role for leisure-time exercise, while prolonged sedentary exposure was associated with greater MetS risk \citep{Kuwahara2016,Browne2017,Huang2017Taiwan,Motuma2020}. The distinction between leisure-time, commuting and occupational activity is important. Physically demanding work is not necessarily equivalent to structured activity that improves cardiorespiratory fitness, and commuting activity was not consistently associated with lower MetS risk. This pattern is consistent with the idea that the metabolic benefit of activity depends not only on total movement, but also on intensity, recovery and the context in which activity occurs.

Dietary evidence also supports workplace-level intervention, although observational findings are more susceptible to measurement error and confounding. The inverse association between Mediterranean-style eating and MetS in a large employee cohort \citep{ReinosoBarbero2026}, together with dietary-pattern findings among health workers \citep{Sharifi2025}, suggests that diet quality is relevant across occupational settings. The randomised canteen trial is especially informative because it demonstrates that modifying the food environment, rather than relying solely on individual advice, can improve multiple metabolic markers \citep{Ma2025}. Findings among night-shift healthcare workers further suggest that meal timing and work schedules may influence metabolic risk in ways that conventional diet-quality measures do not fully capture \citep{dosReis2026}.

The intervention literature is encouraging but heterogeneous. Randomised programmes demonstrated meaningful improvements in MetS severity, waist circumference, blood pressure, glucose regulation and lipid measures \citep{Haufe2019,Ma2025,Bogale2025}. Quasi-experimental studies generally showed the same direction of effect, but their lack of randomisation or concurrent controls reduces causal certainty. The practical implication is not that one programme can be recommended universally, but that metabolic risk appears modifiable when interventions are sufficiently intensive, sustained and adapted to the workplace. Programmes that combine individual support with changes to the work environment may be more promising than approaches that depend on motivation alone.

The diversity of occupations represented in the evidence base is both a strength and a challenge. Office workers, police officers, drivers, industrial workers, healthcare staff, military personnel and university employees experience very different combinations of sitting, shift work, physical demand, psychosocial stress and food access. This heterogeneity helps demonstrate that MetS is relevant across sectors, but it also limits direct comparison and makes a single pooled effect estimate difficult to interpret. Occupational health strategies should therefore be adapted to the risk profile of the workplace rather than assumed to transfer unchanged across sectors.

The review also highlights an important geographical consideration. Several recent studies came from Ghana, Ethiopia and other settings outside the high-income countries that historically dominate occupational-health research. This broadens the relevance of the evidence, but substantial gaps remain. More longitudinal and intervention research is needed in low- and middle-income settings, where occupational health infrastructure, dietary environments and access to preventive care may differ substantially from those in which many workplace programmes were developed.

This review has several strengths. It focuses specifically on employed populations, integrates psychosocial and behavioural determinants within one framework, distinguishes the evidential weight of randomised, prospective and cross-sectional designs, and verifies key effect estimates against source reports. Design-specific JBI appraisal was used to qualify interpretation rather than reduce methodological quality to a single score. Detailed study characteristics and item-level appraisal results are retained in the supplementary material so that the main text can focus on synthesis.

The review also has limitations. Complete historical screening records were unavailable, so a fresh search and re-screening process was required. The search used PubMed/MEDLINE as the primary bibliographic source, supplemented by targeted searching and citation chasing; studies not indexed there may have been missed. Only English-language peer-reviewed studies were eligible. The review was not prospectively registered, no formal publication-bias analysis was performed, and no GRADE certainty framework was applied. Finally, substantial heterogeneity in occupational setting, exposure definition, MetS criteria, intervention design and effect measure precluded an overall meta-analysis.

For occupational-health practice, the evidence favours a multidimensional approach. Psychosocial risk management, healthier food environments, opportunities for physical activity, reduction of prolonged sitting and targeted cardiometabolic screening are likely to be more useful when considered together than when implemented as isolated activities. For research, future studies should use consistent MetS definitions, clearly specify occupational exposures, report intervention fidelity and adherence, and prioritise prospective and pragmatic randomised designs with sufficient follow-up to assess sustained metabolic change.

\section{Conclusion}

Metabolic syndrome among working adults is associated with a combination of occupational and behavioural factors. Prospective evidence supports adverse work stress as a relevant risk marker, while physical activity, sedentary behaviour and diet show broadly consistent associations with metabolic health. Workplace-based lifestyle interventions can improve several MetS components, although the strength of inference varies by design. The workplace is therefore both a setting in which metabolic risk accumulates and a practical setting in which it can be modified.

\section*{Registration and protocol}
This review was not prospectively registered and no prospectively registered protocol was available. The eligibility criteria, search strategy, extraction framework and appraisal procedures are documented in the project repository.

\section*{Funding}
No external funding was received for this systematic review.

\section*{Competing interests}
None declared.

\section*{Ethics approval}
Ethics approval was not required because this systematic review synthesises previously published studies and did not involve collection of new participant-level data.

\section*{Data availability}
Study-level extraction tables, appraisal files, search documentation and reproducibility materials are available at \url{https://github.com/OG-Kyere/workplace-stress-metabolic-syndrome-systematic-review}.

\section*{Author contributions}
Gideon Ofosu Kyere conceived the review update, developed the eligibility framework, conducted the updated search and screening, extracted and synthesised the evidence, performed methodological appraisal, and drafted and revised the manuscript.

\section*{Use of artificial intelligence}
Generative artificial-intelligence tools were used for editorial and organisational assistance during manuscript development. All eligibility decisions, methodological judgments, extracted study results, interpretations and final manuscript content were reviewed by the author against the underlying sources.

\bibliographystyle{unsrtnat}
\bibliography{oem_submission_references}

\end{document}